\chapter{Codetermination}
\label{sec:10}
A floor isn't just a list of forbidden acts — it's an allocation of authority, and there are four distinct powers in it: to \emph{write} the floor, \emph{amend} it, \emph{inspect} whether a deployed system still holds its commitment to it, and \emph{enforce} it by stopping the deployment. They needn't sit with the same institution and usually shouldn't. Calling every arrangement “oversight” hides the differences. The problem begins when one institution holds all four: one that can write, amend, inspect itself, and decide whether a violation has consequences has only described what it currently intends to do.

The powers come apart in practice. Google's AI Principles, published in 2018 after the Maven protests, excluded weapons and surveillance \autocite{castelino2018google,conger2018maven,pichai2018principles} for nearly seven years until the pledge was removed in 2025 \autocite{google2025principles}. The demand became a \emph{policy}: the objectors could hold the company to it in public, and nobody outside held it, so it lasted exactly as long as the company's willingness, and ending it had to be announced. That's the partial case: a transfer held at the institution's pleasure, and the 2025 announcement is what taking one back looks like. 

A refusal in the standard licence holds whichever vendor an operator uses.

The Gebru and Mitchell dismissals are harder, because the power that failed was inspection: the company employed researchers to examine its own systems and dismissed them when the examination became costly \autocite{dickey2020gebru,dickey2021mitchell}.

\textbf{Codetermination.} Codetermination turns protest into a constitutional voice, and it does a second thing the first hides. There is always a surround in which people learn what is acceptable, and the only question is who composes it; for most adults, for most of the day, the employer does \autocite{anderson2017private}. The radicalization literature and the rescue literature are one finding read from both ends: people join because their friends did, and a milieu in which an act has stopped registering as a choice produces the act without needing a bad person \autocite{sageman2008leaderless,mccauley2008mechanisms,malthaner2014radical}. Browning's men shot because the alternative was leaving their comrades to do it; the Oliners' rescuers hid Jews because the people they were attached to made it thinkable. Codetermination and exit are then not only labor rights. They change who composes the surround. I would give part of the four powers to the people who build, label, train and run these systems.

\begin{itemize}
\item \emph{An individual right to refuse.} Every worker may refuse work on any category of use without penalty. The right is backed by anti-retaliation protection that covers hiring, reassignment and contract renewal, and, under capture, by an exit the state doesn't administer through the employer: an offer of work from someone else, drafted in Appendix~\ref{sec:A} on Tcherneva's model, which the offer of work for a bearer in Part VII inherits. The protection is the remedy and the guarantee is the exit. Under capture only the second survives, which is why the floor needs both.
\item \emph{A collective veto, narrowly.} A works council at a lab, contractor or vendor holds a veto only over uses the floor already forbids. It enforces the floor from inside; it doesn't write a second one.
\item \emph{Information, objection and arbitration for everything else.} The council receives information about categories of deployment. Its objection pauses the work, and the dispute goes to arbitration by a body the employer doesn't appoint, on a fixed timetable, so the pause can't become an indefinite block or an ignored letter.
\item \emph{Classified work.} Elected delegates with clearance see the specifics. A denial of clearance to an elected delegate is reviewable by the floor court, and the council may elect alternates, so the state can't choose which workers get to object.
\end{itemize}

Annotators get the same structure: a right to refuse a category they judge wrong, and collective bargaining over the taxonomies they apply, not only a route to file a complaint.

Codetermination gives voice to the people who make these systems, and they are not the people the systems act on. A workforce whose livelihood depends on a contract has interests that point toward it, and the Maven history shows staff objecting at one firm and not at the next. So the voice divides. Over the conditions of their work (what they label, what they are exposed to, what taxonomies they apply), the workers decide. Over the uses the work is put to, the affected public that notice organizes comes first. The organizations representing each published class have standing before the arbitration body, may object to a category of use on their own, and may bring an objection when the works council doesn't.
